% GENERATED FILE. Edit canonical lessons, then run npm run generate:books.
% canonical-source-hash: fnv1a64:da25611d093fc9b0
% canonical-lessons: KA-C365-hodedata, KA-C365-himse, KA-C365-kaidi, KA-C365-nuku, KA-C365-gastu

\chapter{A Fight Breaks Out}
\label{ch:ka-a-fight-breaks-out}
\hlchaptermodality{365}

\begin{chapteropening}
\textbf{By the end of this chapter:} \emph{I can say a fistfight, violence, a prisoner, to shove and a patrol.}

Complete the last lesson of chapter 365: I can say a fistfight, violence, a prisoner, to shove and a patrol.
\end{chapteropening}

\section[\texorpdfstring{\kn{ಹೊಡೆದಾಟ} (hoḍedāṭa)}{ಹೊಡೆದಾಟ (hoḍedāṭa)}]{\kn{ಹೊಡೆದಾಟ} (hoḍedāṭa) — a fistfight; a brawl}

\label{lesson:KA-C365-hodedata}



\begin{quote}
Before the new one: say the Kannada for a preface, then the Kannada for the gist.
\end{quote}

\subsection*{You'll want to know: \kn{ಹೊಡೆದಾಟ}}
\textbf{\kn{ಹೊಡೆದಾಟ}} — \emph{hoḍedāṭa} — \textquotedblleft{}a fistfight; a brawl\textquotedblright{}.

\begin{quote}
\kn{ಪಂದ್ಯದ} \kn{ನಂತರ} \kn{ಎರಡು} \kn{ಗುಂಪುಗಳ} \kn{ನಡುವೆ} \kn{ಹೊಡೆದಾಟ} \kn{ಆಯಿತು}.
\end{quote}

\begin{quote}
\emph{pandyada nantara eraḍu gumpugaḷa naḍuve hoḍedāṭa āyitu.}
\end{quote}

\begin{quote}
After the match there was a fight between two groups.
\end{quote}

\kn{ಹೊಡೆದು}, having hit, and \kn{ಆಟ}, play: hitting back and forth. The verb is \kn{ಹೊಡೆದಾಡು}, to come to blows.

\begin{cousinweb}
\textbf{\kn{ಹೊಡೆದಾಟ}} (\emph{hoḍedāṭa}) is a native Dravidian compound: \kn{ಹೊಡೆದು}, \textquotedblleft{}having struck\textquotedblright{}, from \emph{hoḍe}, \textquotedblleft{}to strike, to beat\textquotedblright{} (DEDR 4252), and \kn{ಆಟ}, \textquotedblleft{}play\textquotedblright{}, from \emph{āḍu} (DEDR 347).
\end{cousinweb}

\begin{grammarlens}[title={What you've built}]
One more word for telling what happened.
\end{grammarlens}

\subsection*{Your turn}
\begin{itemize}
\raggedright
  \item \emph{Say it:} \emph{hoḍedāṭa}
  \item \emph{Say it:} \emph{hoḍedāṭa}, once more
  \item \emph{Say it:} \emph{pandyada nantara eraḍu gumpugaḷa naḍuve hoḍedāṭa āyitu.} — After the match there was a fight between two groups.
  \item \emph{Recall:} say the Kannada for a preface, then the Kannada for the gist, then say \emph{hoḍedāṭa} again
\end{itemize}

\begin{tcolorbox}[breakable,colback=teal!4,colframe=teal!35!black,title={Before you move on}]
What does \kn{ಹೊಡೆದಾಟ} mean? (\textquotedblleft{}A fistfight; a brawl\textquotedblright{}.) Say it once more.
\end{tcolorbox}
\section[\texorpdfstring{\kn{ಹಿಂಸೆ} (hiṃse)}{ಹಿಂಸೆ (hiṃse)}]{\kn{ಹಿಂಸೆ} (hiṃse) — violence; cruelty}

\label{lesson:KA-C365-himse}



\begin{quote}
Before the new one: say the Kannada for the gist, then the Kannada for a fistfight.
\end{quote}

\subsection*{You'll want to know: \kn{ಹಿಂಸೆ}}
\textbf{\kn{ಹಿಂಸೆ}} — \emph{hiṃse} — \textquotedblleft{}violence; cruelty\textquotedblright{}.

\begin{quote}
\kn{ಗಲಭೆಯಲ್ಲಿ} \kn{ತುಂಬಾ} \kn{ಹಿಂಸೆ} \kn{ನಡೆಯಿತು}.
\end{quote}

\begin{quote}
\emph{galabheyalli tumbā hiṃse naḍeyitu.}
\end{quote}

\begin{quote}
There was a lot of violence in the riot.
\end{quote}

\kn{ಹಿಂಸೆ} \kn{ಕೊಡು} is to torment; \kn{ಅಹಿಂಸೆ} is non-violence.

\begin{cousinweb}
\textbf{\kn{ಹಿಂಸೆ}} (\emph{hiṃse}) is borrowed from Sanskrit \emph{hiṃsā}, \textquotedblleft{}injury, harm, hurt\textquotedblright{} (Monier-Williams).
\end{cousinweb}

\begin{grammarlens}[title={What you've built}]
One more word for telling what happened.
\end{grammarlens}

\subsection*{Your turn}
\begin{itemize}
\raggedright
  \item \emph{Say it:} \emph{hiṃse}
  \item \emph{Say it:} \emph{hiṃse}, once more
  \item \emph{Say it:} \emph{galabheyalli tumbā hiṃse naḍeyitu.} — There was a lot of violence in the riot.
  \item \emph{Recall:} say the Kannada for the gist, then the Kannada for a fistfight, then say \emph{hiṃse} again
\end{itemize}

\begin{tcolorbox}[breakable,colback=teal!4,colframe=teal!35!black,title={Before you move on}]
What does \kn{ಹಿಂಸೆ} mean? (\textquotedblleft{}Violence; cruelty\textquotedblright{}.) Say it once more.
\end{tcolorbox}
\section[\texorpdfstring{\kn{ಕೈದಿ} (kaidi)}{ಕೈದಿ (kaidi)}]{\kn{ಕೈದಿ} (kaidi) — a prisoner}

\label{lesson:KA-C365-kaidi}



\begin{quote}
Before the new one: say the Kannada for a fistfight, then the Kannada for violence.
\end{quote}

\subsection*{You'll want to know: \kn{ಕೈದಿ}}
\textbf{\kn{ಕೈದಿ}} — \emph{kaidi} — \textquotedblleft{}a prisoner\textquotedblright{}.

\begin{quote}
\kn{ರಾತ್ರಿ} \kn{ಒಬ್ಬ} \kn{ಕೈದಿ} \kn{ಸೆರೆಮನೆಯಿಂದ} \kn{ತಪ್ಪಿಸಿಕೊಂಡನು}.
\end{quote}

\begin{quote}
\emph{rātri obba kaidi seremaneyinda tappisikoṇḍanu.}
\end{quote}

\begin{quote}
At night a prisoner escaped from the prison.
\end{quote}

Plural: \kn{ಕೈದಿಗಳು}, \emph{kaidigaḷu}.

A prisoner is held in a \kn{ಸೆರೆಮನೆ}, a prison; \kn{ಕೈದಿ} \kn{ಸಂಖ್ಯೆ} is his number.

\begin{cousinweb}
\textbf{\kn{ಕೈದಿ}} (\emph{kaidi}) is a loan from Arabic \emph{qaidī}, \textquotedblleft{}a captive, a prisoner\textquotedblright{}, which came through Persian and Hindustani.
\end{cousinweb}

\begin{grammarlens}[title={What you've built}]
One more word for telling what happened.
\end{grammarlens}

\subsection*{Your turn}
\begin{itemize}
\raggedright
  \item \emph{Say it:} \emph{kaidi}
  \item \emph{Say it:} \emph{kaidi}, once more
  \item \emph{Say it:} \emph{rātri obba kaidi seremaneyinda tappisikoṇḍanu.} — At night a prisoner escaped from the prison.
  \item \emph{Recall:} say the Kannada for a fistfight, then the Kannada for violence, then say \emph{kaidi} again
\end{itemize}

\begin{tcolorbox}[breakable,colback=teal!4,colframe=teal!35!black,title={Before you move on}]
What does \kn{ಕೈದಿ} mean? (\textquotedblleft{}A prisoner\textquotedblright{}.) Say it once more.
\end{tcolorbox}
\section[\texorpdfstring{\kn{ನೂಕು} (nūku)}{ನೂಕು (nūku)}]{\kn{ನೂಕು} (nūku) — to shove; to push}

\label{lesson:KA-C365-nuku}



\begin{quote}
Before the new one: say the Kannada for violence, then the Kannada for a prisoner.
\end{quote}

\subsection*{You'll want to know: \kn{ನೂಕು}}
\textbf{\kn{ನೂಕು}} — \emph{nūku} — \textquotedblleft{}to shove; to push\textquotedblright{}.

\begin{quote}
\kn{ಜನಸಂದಣಿಯಲ್ಲಿ} \kn{ಯಾರೋ} \kn{ನನ್ನನ್ನು} \kn{ನೂಕಿದರು}.
\end{quote}

\begin{quote}
\emph{janasandaṇiyalli yārō nannannu nūkidaru.}
\end{quote}

\begin{quote}
Someone shoved me in the crowd.
\end{quote}

Past (I …): \kn{ನೂಕಿದೆ}, \emph{nūkide}.

Who is pushed takes -\kn{ಅನ್ನು}. \kn{ನೂಕುನುಗ್ಗಲು} is a crush of people pushing.

\begin{cousinweb}
\textbf{\kn{ನೂಕು}} (\emph{nūku}) is native Dravidian: DEDR 3722 lists Kannada \emph{nūku}, \textquotedblleft{}to shove, to push, to thrust aside\textquotedblright{}, beside Telugu \emph{nūku}.
\end{cousinweb}

\begin{grammarlens}[title={What you've built}]
One more word for telling what happened.
\end{grammarlens}

\subsection*{Your turn}
\begin{itemize}
\raggedright
  \item \emph{Say it:} \emph{nūku}
  \item \emph{Say it:} \emph{nūku}, once more
  \item \emph{Say it:} \emph{janasandaṇiyalli yārō nannannu nūkidaru.} — Someone shoved me in the crowd.
  \item \emph{Recall:} say the Kannada for violence, then the Kannada for a prisoner, then say \emph{nūku} again
\end{itemize}

\begin{tcolorbox}[breakable,colback=teal!4,colframe=teal!35!black,title={Before you move on}]
What does \kn{ನೂಕು} mean? (\textquotedblleft{}To shove; to push\textquotedblright{}.) Say it once more.
\end{tcolorbox}
\section[\texorpdfstring{\kn{ಗಸ್ತು} (gastu)}{ಗಸ್ತು (gastu)}]{\kn{ಗಸ್ತು} (gastu) — a patrol; a beat}

\label{lesson:KA-C365-gastu}



\begin{quote}
Before the new one: say the Kannada for a prisoner, then the Kannada for to shove.
\end{quote}

\subsection*{You'll want to know: \kn{ಗಸ್ತು}}
\textbf{\kn{ಗಸ್ತು}} — \emph{gastu} — \textquotedblleft{}a patrol; a beat\textquotedblright{}.

\begin{quote}
\kn{ರಾತ್ರಿ} \kn{ಪೊಲೀಸರು} \kn{ಬೀದಿಗಳಲ್ಲಿ} \kn{ಗಸ್ತು} \kn{ತಿರುಗುತ್ತಾರೆ}.
\end{quote}

\begin{quote}
\emph{rātri polīsaru bīdigaḷalli gastu tiruguttāre.}
\end{quote}

\begin{quote}
At night the police patrol the streets.
\end{quote}

Police on their round are \kn{ಗಸ್ತಿನಲ್ಲಿ}, on patrol.

\begin{cousinweb}
\textbf{\kn{ಗಸ್ತು}} (\emph{gastu}) is a loan from Persian \emph{gašt}, \textquotedblleft{}a walk, a round, a patrol\textquotedblright{}, which came through Hindustani.
\end{cousinweb}

\begin{grammarlens}[title={What you've built}]
One more word for telling what happened.
\end{grammarlens}

\subsection*{Your turn}
\begin{itemize}
\raggedright
  \item \emph{Say it:} \emph{gastu}
  \item \emph{Say it:} \emph{gastu}, once more
  \item \emph{Say it:} \emph{rātri polīsaru bīdigaḷalli gastu tiruguttāre.} — At night the police patrol the streets.
  \item \emph{Recall:} say the Kannada for a prisoner, then the Kannada for to shove, then say \emph{gastu} again
\end{itemize}

\begin{tcolorbox}[breakable,colback=teal!4,colframe=teal!35!black,title={Before you move on}]
What does \kn{ಗಸ್ತು} mean? (\textquotedblleft{}A patrol; a beat\textquotedblright{}.) Say it once more.
\end{tcolorbox}
